\section{The algorithm and its bit complexity}\label{sec:algorithm}

The preceding sections prove existence of a successful choice at every
stage. During pairing, the algorithm chooses a minimizer of the discounted
accounts; during repair, it enumerates a polynomial family of simultaneous
switches and accepts the first graph passing an exact spectral test.
We now verify that these operations fit the stated bit model, and then
assemble the theorem. All computations use rational
numbers, except for real quantities used solely to analyze their bounds.


\begin{proposition}\label{prop:bit-complexity}
For each fixed $d$ and the fixed choices in Section~\ref{subsec:hierarchy}, every
stage of the preceding construction is implementable with polynomially
many bit operations in $n$.
\end{proposition}

\begin{proof}
\textit{Rational parameters.}
Choose dyadic rational representatives for all power scales, within the
fixed factors allowed above.  Comparisons for a fixed rational exponent
can be made by integer powers, so these representatives require
$O_d(\log n)$ bits and are computable in polynomial time.  Integer power
comparisons similarly determine the logarithmic radii.  Approximate the
primary $h$ values to sufficiently high inverse-polynomial precision,
strictly below $1/\sqrt q$; choose the small auxiliary spacings rationally
and exactly equally spaced where required.  The strict exponent margins
preserve every estimate.  The piecewise-polynomial cutoff for the smooth
reference can have rational breakpoints and coefficients and fixed
finite smoothness, for example by rescaling the integral of
$t^j(1-t)^j$ at its transition intervals.  Its Fourier transform is used
only in the proof of an estimate, never by the algorithm.

\textit{Candidate computations.}
There are at most $O_d(n^2)$ choices at a pairing step and $O_d(n)$ steps.
Distances, stub eligibility, and exceptional-event flags are finite graph
computations.  The precision matrices have rational entries of
polynomial bit length.  Their inverses on $\mathbf1^\perp$ can be
computed in a rational contrast basis or after adding the constant
projection.  The coefficients of the path series are computed from the
displayed matrix recurrence, of length $O_d(\log n)$.  The frozen
reference is evaluated from its explicit Woodbury formula at the current
matching, rather than through an uncontrolled sequence of rational
expansions.

The structural completion estimators are computed with the colorful
dynamic program described in the structural section.  Its shapes have
$O_d(\log n)$ vertices and bounded treewidth once the fixed excess cutoff
has been chosen.  The factor $2^{O(k)}$ for the color subset and the
$\exp(k)k^{O(\log k)}\log n$ hash-family size are polynomial in $n$ when
$k=O_d(\log n)$.  Pending-edge probabilities are rational products with
$O_d(\log n)$ factors.  These observations give polynomial time and bit
length for every such estimator.

All moments have fixed integer order.  Traces, squares, even powers,
positive parts of rational numbers, comparator values, and uniform
conditional expectations are therefore rational and computable in
polynomial time.  There are polynomially many accounts, including the
fixed contrast labels and their pairs.

\textit{Discounts.}
At a step the expectation ratio for an account is computed from its
undiscounted local statistic, its fixed phase baseline, and its
comparator and event flags.  It is not computed recursively from the
discounted total.  A sum over the polynomially many candidates has
polynomial bit length.  A discount is a product of polynomially many
such ratios; consequently it too has polynomial bit length.  Dynamic
comparators are explicit rational functions or products of polynomially
many polynomial-bit factors.  Sums and comparisons of the discounted
accounts can thus be carried out exactly.

For completeness, the polynomial bit bounds for static matrix inverses
follow from the determinant formula: after a common denominator is
cleared, minors of a polynomial-bit matrix of dimension $O_d(n)$ have
polynomial bit length by the determinant bound.  Exact rational symmetric
inertia can be computed by a nonzero diagonal pivot, or, if all diagonals
vanish and a nonzero off-diagonal remains, by an invertible $2$-by-$2$
pivot.  The remaining zero block contributes zero eigenvalues.
Fraction-free elimination \cite[Section~I]{Bareiss1968} or the minor formulas
give the same bit bound
for intermediate quantities.

\textit{Repair and verification.}
The repair section constructs inverse-polynomially separated cutoffs,
rational polynomial approximations to the required spectral functions,
and a polynomial-size prime-field sample space.  Enumerating that sample
space and performing the graph updates therefore takes polynomial time.
The precision exponent is fixed after the analytic margins have been
chosen.  In particular, after affine rescaling to $[0,1]$, if a spectral
cutoff function has a known Lipschitz bound $L_f\le n^c$ and the target
error is $n^{-C}$, its
Bernstein approximation may be taken of degree
$D=\lceil n^{2(c+C)}\rceil$, polynomial in $n$ for fixed $c,C$.
Its rational sample values and polynomial coefficients have polynomial
bit length.  Increasing the fixed precision and prime-denominator
exponents pays for both the tilt factor $n/J_T$ and the total
probability-rounding error.

For each simple regular candidate, take the explicit rational contrast basis
\[
 B=[e_1-e_n,\ldots,e_{n-1}-e_n].
\]
Its columns are linearly independent and span $\mathbf1^\perp$.
The condition
\[
 B^T(4qI-A_G^2)B
 =4qB^TB-(A_GB)^T(A_GB)\succ0
\]
is equivalent to the strict nontrivial spectral bound and is decidable
by exact rational inertia.  Indeed, this congruence tests the quadratic
form on all of $\mathbf1^\perp$; an orthonormal basis is unnecessary.
Regularity makes that subspace invariant under $A_G$, so the form is
positive exactly when every eigenvalue there has square less than $4q$.
The repair theorem guarantees that the search
finds a passing candidate with a strict spectral inequality.

All constants, polynomial degrees, moment orders, and approximation
exponents depend only on $d$.  A sufficiently large finite exponent
$k_d$ bounds their combined bit cost.  This bound holds once the finitely many large-$n$ requirements are met.
\end{proof}

\begin{proof}[Proof of Theorem~\ref{thm:main}]
Start with $d$ distinguishable stubs at each of $n$ vertices.  Apply
Proposition~\ref{prop:early-phase}, rejecting stub
pairs whose vertices have distance at most the girth exclusion radius.
The rejected-pair graph has maximum degree $O_d(q^{g_0})=o(l_0)$, so the
candidate set is nonempty throughout this phase.  Apply
Proposition~\ref{prop:cleanup} to reach $l_0$ unresolved stubs.  The structural estimates guarantee that
these stubs lie at distinct vertices and are mutually separated, with
no nonempty short nonbacktracking return at any remaining vertex.

The reference construction and Proposition~\ref{prop:boundary-completion} then complete the
configuration to a simple $d$-regular graph.  For every primary parameter
and both signs, its precision is positive; at the closest parameter this
gives
\[
 \|A_G|_{\mathbf1^\perp}\|<s+e/20.
\]
The final entry comparison and the diagonal derivative bound for the same
frozen reference give the local spectral mass estimate by
Lemma~\ref{lem:repair-projector}.  The retained path estimates at the
separation stop, together with that entry comparison, give the typical
bulk edge blocks by Lemma~\ref{lem:repair-bulk}.  These are precisely the
remaining hypotheses of Theorem~\ref{thm:exact-repair}.  Apply that theorem.  This produces a simple
$d$-regular graph on the same $n$ vertices with every nonconstant
eigenvalue strictly in $(-s,s)$.

The exact-repair theorem also proves connectedness and nonbipartiteness.
Proposition~\ref{prop:bit-complexity} supplies a finite exponent $k_d$
for the bit cost. Choose $n_0(d)$ to satisfy the finitely many
large-$n$ requirements in the construction. This proves all parts of
Theorem~\ref{thm:main}.
\end{proof}
